% GENERATED FILE. Edit canonical lessons, then run npm run generate:books.
% canonical-source-hash: fnv1a64:0a074eec3316b177
% canonical-lessons: HI-C262-barhai, HI-C262-mochi, HI-C262-pujari, HI-C262-gahna, HI-C262-kangan

\chapter{A carpenter, A cobbler, A priest, A piece of jewellery, A bangle}
\label{ch:hi-a-carpenter-a-cobbler-a-priest-a-piece-of-jewellery-a-bangle}
\hlchaptermodality{262}

\begin{chapteropening}
\textbf{By the end of this chapter:} \emph{I can say a carpenter, a cobbler, a priest, a piece of jewellery and a bangle.}

Complete the last lesson of chapter 262: I can say a carpenter, a cobbler, a priest, a piece of jewellery and a bangle.
\end{chapteropening}

\section[\texorpdfstring{\hi{बढ़ई} (baṛhaī)}{बढ़ई (baṛhaī)}]{\hi{बढ़ई} (baṛhaī) — a carpenter}

\label{lesson:HI-C262-barhai}



\begin{quote}
Before the new one: say the Hindi for a blacksmith, then the Hindi for a goldsmith.
\end{quote}

\subsection*{You'll want to know: \hi{बढ़ई}}
\textbf{\hi{बढ़ई}} — \emph{baṛhaī} — \textquotedblleft{}a carpenter\textquotedblright{}.

\begin{cousinweb}
\textbf{\hi{बढ़ई}} (\emph{baṛhaī}) is a tadbhava, a word inherited from Sanskrit and reshaped by sound change. It comes from \emph{vardhaki}, \textquotedblleft{}carpenter\textquotedblright{} (CDIAL 11375).
\end{cousinweb}

\begin{grammarlens}[title={What you've built}]
One more word for this chapter.
\end{grammarlens}

\subsection*{Your turn}
\begin{itemize}
\raggedright
  \item \emph{Say it:} \emph{baṛhaī}
  \item \emph{Say it:} \emph{baṛhaī}, once more
  \item \emph{Say it:} say \emph{baṛhaī}
  \item \emph{Recall:} say the Hindi for a blacksmith, then the Hindi for a goldsmith, then say \emph{baṛhaī} again
\end{itemize}

\begin{tcolorbox}[breakable,colback=teal!4,colframe=teal!35!black,title={Before you move on}]
What does \hi{बढ़ई} mean? (\textquotedblleft{}A carpenter\textquotedblright{}.) Say it once more.
\end{tcolorbox}
\section[\texorpdfstring{\hi{मोची} (mochī)}{मोची (mochī)}]{\hi{मोची} (mochī) — a cobbler}

\label{lesson:HI-C262-mochi}



\begin{quote}
Before the new one: say the Hindi for a goldsmith, then the Hindi for a carpenter.
\end{quote}

\subsection*{You'll want to know: \hi{मोची}}
\textbf{\hi{मोची}} — \emph{mochī} — \textquotedblleft{}a cobbler\textquotedblright{}.

\begin{cousinweb}
\textbf{\hi{मोची}} (\emph{mochī}) goes back to Sanskrit \emph{mocika}, \textquotedblleft{}tanner and cobbler\textquotedblright{}, recorded in Buddhist Sanskrit and built on \emph{moca}, \textquotedblleft{}shoe\textquotedblright{} (CDIAL 10349). CDIAL takes \emph{moca} from an Iranian source, comparing Pahlavi (Middle Persian) \emph{mōč}, \textquotedblleft{}shoe\textquotedblright{}.
\end{cousinweb}

\begin{grammarlens}[title={What you've built}]
One more word for this chapter.
\end{grammarlens}

\subsection*{Your turn}
\begin{itemize}
\raggedright
  \item \emph{Say it:} \emph{mochī}
  \item \emph{Say it:} \emph{mochī}, once more
  \item \emph{Say it:} say \emph{mochī}
  \item \emph{Recall:} say the Hindi for a goldsmith, then the Hindi for a carpenter, then say \emph{mochī} again
\end{itemize}

\begin{tcolorbox}[breakable,colback=teal!4,colframe=teal!35!black,title={Before you move on}]
What does \hi{मोची} mean? (\textquotedblleft{}A cobbler\textquotedblright{}.) Say it once more.
\end{tcolorbox}
\section[\texorpdfstring{\hi{पुजारी} (pujārī)}{पुजारी (pujārī)}]{\hi{पुजारी} (pujārī) — a priest}

\label{lesson:HI-C262-pujari}



\begin{quote}
Before the new one: say the Hindi for a carpenter, then the Hindi for a cobbler.
\end{quote}

\subsection*{You'll want to know: \hi{पुजारी}}
\textbf{\hi{पुजारी}} — \emph{pujārī} — \textquotedblleft{}a priest\textquotedblright{}.

\begin{cousinweb}
\textbf{\hi{पुजारी}} (\emph{pujārī}) is \emph{pūjā}, \textquotedblleft{}worship\textquotedblright{}, from Sanskrit \emph{pūjā} (CDIAL 8317), with the ending \emph{-ārī}, \textquotedblleft{}one who does it\textquotedblright{}.
\end{cousinweb}

\begin{grammarlens}[title={What you've built}]
One more word for this chapter.
\end{grammarlens}

\subsection*{Your turn}
\begin{itemize}
\raggedright
  \item \emph{Say it:} \emph{pujārī}
  \item \emph{Say it:} \emph{pujārī}, once more
  \item \emph{Say it:} say \emph{pujārī}
  \item \emph{Recall:} say the Hindi for a carpenter, then the Hindi for a cobbler, then say \emph{pujārī} again
\end{itemize}

\begin{tcolorbox}[breakable,colback=teal!4,colframe=teal!35!black,title={Before you move on}]
What does \hi{पुजारी} mean? (\textquotedblleft{}A priest\textquotedblright{}.) Say it once more.
\end{tcolorbox}
\section[\texorpdfstring{\hi{गहना} (gahnā)}{गहना (gahnā)}]{\hi{गहना} (gahnā) — a piece of jewellery}

\label{lesson:HI-C262-gahna}



\begin{quote}
Before the new one: say the Hindi for a cobbler, then the Hindi for a priest.
\end{quote}

\subsection*{You'll want to know: \hi{गहना}}
\textbf{\hi{गहना}} — \emph{gahnā} — \textquotedblleft{}a piece of jewellery\textquotedblright{}.

\begin{cousinweb}
\textbf{\hi{गहना}} (\emph{gahnā}) is a tadbhava, a word inherited from Sanskrit and reshaped by sound change. It comes from \emph{grahaṇa}, \textquotedblleft{}seizing; putting on (clothes)\textquotedblright{} (CDIAL 4364). CDIAL glosses the Hindi word \textquotedblleft{}jewel\textquotedblright{}: a thing put on.
\end{cousinweb}

\begin{grammarlens}[title={What you've built}]
One more word for this chapter.
\end{grammarlens}

\subsection*{Your turn}
\begin{itemize}
\raggedright
  \item \emph{Say it:} \emph{gahnā}
  \item \emph{Say it:} \emph{gahnā}, once more
  \item \emph{Say it:} say \emph{gahnā}
  \item \emph{Recall:} say the Hindi for a cobbler, then the Hindi for a priest, then say \emph{gahnā} again
\end{itemize}

\begin{tcolorbox}[breakable,colback=teal!4,colframe=teal!35!black,title={Before you move on}]
What does \hi{गहना} mean? (\textquotedblleft{}A piece of jewellery\textquotedblright{}.) Say it once more.
\end{tcolorbox}
\section[\texorpdfstring{\hi{कंगन} (kaṅgan)}{कंगन (kaṅgan)}]{\hi{कंगन} (kaṅgan) — a bangle, a bracelet}

\label{lesson:HI-C262-kangan}



\begin{quote}
Before the new one: say the Hindi for a priest, then the Hindi for a piece of jewellery.
\end{quote}

\subsection*{You'll want to know: \hi{कंगन}}
\textbf{\hi{कंगन}} — \emph{kaṅgan} — \textquotedblleft{}a bangle, a bracelet\textquotedblright{}.

\begin{cousinweb}
\textbf{\hi{कंगन}} (\emph{kaṅgan}) goes back to Sanskrit \emph{kaṅkaṇa}, \textquotedblleft{}bracelet\textquotedblright{} (CDIAL 2597). CDIAL lists Hindi \emph{kaknā} and \emph{kaknī} under it, but not \emph{kaṅgan} itself.
\end{cousinweb}

\begin{grammarlens}[title={What you've built}]
One more word for this chapter.
\end{grammarlens}

\subsection*{Your turn}
\begin{itemize}
\raggedright
  \item \emph{Say it:} \emph{kaṅgan}
  \item \emph{Say it:} \emph{kaṅgan}, once more
  \item \emph{Say it:} say \emph{kaṅgan}
  \item \emph{Recall:} say the Hindi for a priest, then the Hindi for a piece of jewellery, then say \emph{kaṅgan} again
\end{itemize}

\begin{tcolorbox}[breakable,colback=teal!4,colframe=teal!35!black,title={Before you move on}]
What does \hi{कंगन} mean? (\textquotedblleft{}A bangle, a bracelet\textquotedblright{}.) Say it once more.
\end{tcolorbox}
